\section{Count-only reconstruction from a fixed local bracket}
\label{sec:v7-bracket}
The preceding acquisition theorem presupposes a bracket whose width
shrinks with the desired accuracy. We now acquire that bracket using
counts, starting from a fixed local interval. A deterministic cap on
all later waiting times prevents an unsuccessful search from producing
an infinite preparation cost. The original locally supplied-bracket
theorem remains useful, with its sharper description of each stage.

Fix a sufficiently small compact physical neighborhood in
Theorem~\ref{thm:v6-three}. The constants of the forward law on this
neighborhood, including a positive lower bound, are regarded as design
information. This is local reconstruction in a specified family, not
a procedure requiring no prior geometric information. For $j=1$ there
are known constants $c_*>0$, $d_*>0$ such that
\begin{equation}\label{eq:v7-search-onset}
 \Pr(N_{g+d}\ge2)\ge c_*d^2\quad(0<d\le d_*),
 \qquad \Pr(N_t\ge2)=0\quad(t\le g).
\end{equation}
These follow from Theorem~\ref{thm:g-stability} and compactness.

\begin{lemma}[Count acquisition of a coarse bracket]
\label{lem:v7-bracket}
Suppose $g\in[L_0,U_0]$, with fixed width $0<W=U_0-L_0\le d_*$.
For $0<h<W$ and $0<\delta<1/2$, a binary count procedure returns
$g_0$ satisfying $|g-g_0|\le h/4$ with probability at least
$1-\delta$. On every history it uses at most
\begin{equation}\label{eq:v7-bracket-cost}
 C h^{-2}\log\left(\frac{2+\log(W/h)}{\delta}\right)
\end{equation}
independent preparations. Its queries use only the indicator of two
or more collisions at programmed times inside $[L_0,U_0]$.
\end{lemma}
\begin{proof}
Maintain an interval $[L,U]$, initially $[L_0,U_0]$, and stop when
its width $w$ is at most $h/2$. At a nonterminal stage query
$t=(L+U)/2$ on
\[
 n(w)=\left\lceil\frac{16}{c_*w^2}
                  \log\frac{M+1}{\delta}\right\rceil
\]
fresh preparations, where
$M=\lceil\log(2W/h)/\log(4/3)\rceil$.
If at least one query succeeds, replace $U$ by $t$; otherwise replace
$L$ by $t-w/4$. In either case the interval is nested in the previous
one and its width is either $w/2$ or $3w/4$. There are at most $M$
stages, regardless of the observations.

Conditional on an interval containing $g$, a success certifies $g<t$
by \eqref{eq:v7-search-onset}. A failure branch can discard $g$ only
if $g<t-w/4$. In this case $t-g>w/4$ and $t-g\le W\le d_*$, so the
probability of seeing no success is at most
$\exp(-n(w)c_*w^2/16)\le\delta/(M+1)$. A union bound over the first
possible incorrect discard proves that all retained intervals contain
$g$ with probability at least $1-\delta$. The midpoint of the final
interval then has error at most $h/4$.

For cost, each nonterminal width exceeds $h/2$ and contracts by at
least a factor $3/4$. Summing reciprocal squared widths backwards
from the last such stage gives $\sum w^{-2}\le C h^{-2}$.
The ceiling contributions add at most $M$, also bounded by
$C h^{-2}$ on the stated range, with a constant depending on fixed
$W$. This proves \eqref{eq:v7-bracket-cost} and its all-history form.
There is no waiting for an event of unknown positive probability in
this search.
\end{proof}

\begin{theorem}[Fully charged local count reconstruction]
\label{thm:v7-fixed-bracket}
Fix $m\ge1$, a compact physical neighborhood as above, and a fixed
interval $[L_0,U_0]$ containing its possible gaps and satisfying the
hypotheses of Lemma~\ref{lem:v7-bracket}. There exist $h_m>0$ and
$C_m<\infty$ such that, for $0<h\le h_m$ and $0<\eta<1/2$, a
count-only procedure estimates the gap, area, and unordered curvature
triple with
\begin{equation}\label{eq:v7-fixed-bracket-errors}
 |\widehat g-g|\le C_mh^{m+1},\qquad
 |\widehat A-A|\le C_mh^m,\qquad
 \operatorname{dist}_{\rm match}(\widehat\kappa,\kappa)
                                      \le C_mh^{m/3}
\end{equation}
with probability at least $1-\eta$. Its total preparation count is
bounded on every history by
\begin{equation}\label{eq:v7-fixed-bracket-cost}
 C_m h^{-(2m+2)}\log(C_m/\eta).
\end{equation}
The procedure uses only binary collision-count reports with flight
numbers $1,2,3$, independent full-phase preparations, and noiseless
elapsed-time recording. In particular curvature accuracy $\varepsilon$
has sufficient total cost
$C_m\varepsilon^{-(6+6/m)}\log(C_m/\eta)$ from a bracket independent
of $\varepsilon$. The supplied information includes the fixed local
family and its uniform bounds, but no accuracy-dependent gap bracket,
contact positions, or separately measured area.
\end{theorem}
\begin{proof}
Use Lemma~\ref{lem:v7-bracket} with failure budget $\eta/4$ to acquire
$g_0$. Conditional on its success, the hypothesis
$|g-g_0|\le h/4$ of Theorem~\ref{thm:v6-count-only} holds.
All subsequent preparations are fresh. Run that theorem's pilot and
amplitude stages with $J=3$, and allocate failure probability
$\eta/2$ to their estimation bounds. Their constants are uniform over
all admissible $g_0$, so conditioning on the entire first stage is
legitimate. On these good events their estimates satisfy
\eqref{eq:v7-fixed-bracket-errors}.

There is an additional issue on a failed first stage: a programmed
pilot window could lie below onset. We therefore modify each pilot
waiting time as follows. If its target
$r=\lceil C_mh^{-2m}\log(C_m/\eta)\rceil$ successes have not occurred
within
\[
 B_m=\left\lceil C_m h^{-(2m+2)}\log(C_m/\eta)\right\rceil
\]
preparations, stop the entire procedure and return a fixed admissible
parameter as a fallback. On first-stage success, every pilot
probability is at least $c_mh^2$, uniformly over its programmed nodes.
Choose the constant in $B_m$ so that $B_mc_mh^2\ge2r$ and
$B_mc_mh^2\ge8\log(4(m+1)/\eta)$. The exponential Markov bound for a
binomial variable gives
\[
 \Pr\{T_r>B_m\}
  =\Pr\{\operatorname{Bin}(B_m,p)<r\}
  \le e^{-B_mp/8}\le\eta/[4(m+1)].
\]
A union bound over the $m+1$ pilots charges at most $\eta/4$ for a
cap on a good first stage. Before a cap the procedure agrees exactly
with the original one, so this modification does not condition or
bias the claimed estimation guarantee. The final amplitude stage has
a predetermined number of preparations, whether or not its estimates
are good. Its compact minimizer is selected by
Lemma~\ref{lem:v7-selection}.

Every history now has a deterministic preparation bound: the first
stage has \eqref{eq:v7-bracket-cost}, each of the fixed number of
pilots has cap $B_m$, and the amplitude stage has the deterministic
bound in the proof of Theorem~\ref{thm:v6-count-only}. For fixed
$m\ge1$, the function
$h^{2m}\log(2+\log(W/h))$ is bounded for $0<h\le h_m<W$.
Thus the first-stage cost is absorbed into
\eqref{eq:v7-fixed-bracket-cost}, uniformly in $\eta<1/2$.
The three allocated failure probabilities sum to at most $\eta$.
Finally choose $h=c_m\varepsilon^{3/m}$, with fixed $c_m$ small enough
to make the curvature bound at most $\varepsilon$.
\end{proof}

The same initial search and cap construction applies to the independent-
normalizer four-amplitude experiment of Theorem~\ref{thm:v6-count-only}
on any specified compact physical family satisfying its hypotheses;
then $J=4$. The search uses only its positive one-flight onset bound,
and no area constraint is needed for that stage.

The first-stage windows are allowed to have offsets larger than $h$.
Accordingly Theorem~\ref{thm:v7-fixed-bracket} is an upper bound for a
larger design class than the all-history shrinking-collar class in
\eqref{eq:v7-cost-lower}. We do not infer global minimax optimality
for this enlarged class from that lower bound. The weighted-exposure
inequality \eqref{eq:v7-kl-exposure} remains applicable whenever its
own query hypotheses hold. The selected-position acquisition theorem
has different data and is unchanged by this count-only construction.
